\chapter{Implementation and Experiments}
\label{ch:implementation}

This chapter presents the Julia implementations of the three power flow modules alongside their numerical validation and performance evaluation. The chapter opens with the definition of a canonical 3-bus test network that serves as a common validation fixture throughout. The DC power flow, AC power flow, and LOPF implementations are then described in turn: each section presents the core source code, discusses the engineering decisions and non-obvious challenges encountered during porting, and reports numerical results compared against the reference PyPSA implementation. The chapter concludes with a systematic performance benchmark spanning networks from 3 to 2{,}000 buses, followed by an analysis that accounts for the observed speedup patterns --- including the opposing trends exhibited by DC power flow and LOPF as problem size grows.

% ──────────────────────────────────────────────────────────────────────────────
\section{Test Network Definition}
\label{sec:test_network_definition}
% ──────────────────────────────────────────────────────────────────────────────

All three implementations are validated on the same standard 3-bus test network, ensuring consistency across comparisons with PyPSA. The network is defined as follows:

\begin{itemize}
    \item \textbf{Bus 1:} Generator G1, capacity $P_1^\text{max} = 400\,\text{MW}$, marginal cost $c_1 = 20\,\text{\euro/MWh}$, reference bus ($\theta_1 = 0$).
    \item \textbf{Bus 2:} Generator G2, capacity $P_2^\text{max} = 300\,\text{MW}$, marginal cost $c_2 = 50\,\text{\euro/MWh}$; load $L_2 = 200\,\text{MW}$.
    \item \textbf{Bus 3:} Load $L_3 = 300\,\text{MW}$, no local generation.
\end{itemize}

All three transmission lines have reactance $x = 0.1\,\text{p.u.}$ (base: $S_\text{base} = 100\,\text{MVA}$, $V_\text{base} = 380\,\text{kV}$), giving a susceptance of:
\begin{equation}
    b = \frac{S_\text{base}}{x} = \frac{100\,\text{MVA}}{0.1\,\text{p.u.}} = 1000\,\text{MW/rad}.
    \label{eq:susceptance_value}
\end{equation}

The susceptance matrix of this network is:
\begin{equation}
    \mathbf{B} = \begin{pmatrix} 2000 & -1000 & -1000 \\ -1000 & 2000 & -1000 \\ -1000 & -1000 & 2000 \end{pmatrix} \text{ MW/rad}.
\end{equation}

% ──────────────────────────────────────────────────────────────────────────────
\section{DC Power Flow Implementation}
% ──────────────────────────────────────────────────────────────────────────────

\subsection{Core Implementation}

The DC power flow solver is implemented in \texttt{julia/dc\_power\_flow.jl}. The core function assembles the susceptance matrix and solves the linear system using Julia's backslash operator, which dispatches to LAPACK for dense matrices:

\begin{lstlisting}[language=Julia, caption={DC power flow: susceptance matrix and linear solve}, label={lst:dc_pf_core}]
function dc_pf_solve(n_buses, lines, generators, loads; baseMVA=100.0)
    # Build susceptance matrix B (n x n)
    B = zeros(n_buses, n_buses)
    for (from, to, x) in lines
        b = baseMVA / x          # b [MW/rad] = S_base / x_pu
        B[from, from] += b
        B[to,   to  ] += b
        B[from, to  ] -= b
        B[to,   from] -= b
    end

    # Net injection vector P = P_gen - P_load
    P = zeros(n_buses)
    for (bus, p) in generators;  P[bus] += p;  end
    for (bus, p) in loads;       P[bus] -= p;  end

    # Remove slack bus (bus 1): solve reduced (n-1) x (n-1) system
    idx   = 2:n_buses
    theta = zeros(n_buses)
    theta[idx] = B[idx, idx] \ P[idx]    # Julia backslash -> LAPACK
    return theta
end
\end{lstlisting}

\subsection{Validation}

For the DC power flow validation, the generator dispatch is fixed at $G_1 = 400\,\text{MW}$ (slack, absorbs imbalance), $G_2 = 100\,\text{MW}$ (the economically optimal dispatch from LOPF Scenario~A). The net injections are:
\begin{equation}
    \mathbf{P} = \begin{pmatrix} 400 \\ 100 - 200 \\ 0 - 300 \end{pmatrix} = \begin{pmatrix} 400 \\ -100 \\ -300 \end{pmatrix}\,\text{MW}.
\end{equation}

Solving the reduced $2 \times 2$ system $\hat{\mathbf{B}}\hat{\boldsymbol{\theta}} = \hat{\mathbf{P}}$ analytically:
\begin{equation}
    \begin{pmatrix} 2000 & -1000 \\ -1000 & 2000 \end{pmatrix}
    \begin{pmatrix} \theta_2 \\ \theta_3 \end{pmatrix}
    = \begin{pmatrix} -100 \\ -300 \end{pmatrix}
    \implies
    \begin{pmatrix} \theta_2 \\ \theta_3 \end{pmatrix}
    = \begin{pmatrix} -1/6 \\ -7/30 \end{pmatrix}\,\text{rad}.
\end{equation}

Table~\ref{tab:dc_validation} compares the Julia and PyPSA results on this network.

\begin{table}[H]
\centering
\caption{DC power flow validation: Julia vs PyPSA (3-bus network, $G_1 = 400\,\text{MW}$, $G_2 = 100\,\text{MW}$). Branch flows agree exactly; voltage angles differ by a constant per-unit convention factor of $V_\text{nom}^2/S_\text{base} = 380^2/100 = 1444$ (see text).}
\label{tab:dc_validation}
\begin{tabular}{lrrl}
\toprule
\textbf{Quantity} & \textbf{Julia} & \textbf{PyPSA} & \textbf{Relationship} \\
\midrule
$\theta_2$ (rad)   & $-$0.1667 & $-1.15\times10^{-4}$ & ratio $= 1444$ \\
$\theta_3$ (rad)   & $-$0.2333 & $-1.62\times10^{-4}$ & ratio $= 1444$ \\
$p_{12}$ (MW)      &  166.7    &  166.7    & $\Delta < 10^{-6}$ \\
$p_{13}$ (MW)      &  233.3    &  233.3    & $\Delta < 10^{-6}$ \\
$p_{23}$ (MW)      &   66.7    &   66.7    & $\Delta < 10^{-6}$ \\
\bottomrule
\end{tabular}
\end{table}

The branch flows --- the physically and economically meaningful output of the DC power flow --- agree to better than $10^{-6}\,\text{MW}$. The voltage angles differ by an exact constant factor of $1444$ because the two implementations adopt different per-unit conventions for the line reactance. PyPSA interprets the \texttt{Line.x} attribute as an \emph{ohmic} value and converts it internally to per-unit by dividing by $V_\text{nom}^2/S_\text{base} = 380^2/100 = 1444\,\Omega$, producing angles of order $10^{-4}$\,rad. The present implementation, following the MATPOWER convention, treats \texttt{x} as already in per-unit and forms the susceptance as $b = S_\text{base}/x$ in MW/rad, producing angles of order $10^{-1}$\,rad. Because the angle scale cancels exactly in the branch-flow expression $p_{km} = b_{km}(\theta_k - \theta_m)$, the resulting flows are identical. The difference is therefore a documented input-convention choice, not a modelling error; the same flows, dispatch, and prices are obtained under both conventions. Notably, this scale difference appears only in the synthetic toy network, where a bare per-unit reactance ($x = 0.1$) is supplied against a $380\,\text{kV}$ base; on the IEEE 14-bus case of Section~\ref{sec:ieee14}, where the MATPOWER data carry physically consistent per-unit reactances, the Julia and PyPSA angles agree to better than $0.001^\circ$ across all buses (Table~\ref{tab:ieee14_dc}). Figure~\ref{fig:dc_pf_result} shows the voltage angles at each bus.

\begin{figure}[H]
    \centering
    \includegraphics[width=\textwidth]{results_validation.png}
    \caption{Validation results on the 3-bus test network (left to right):
    (a) DC Power Flow voltage angles at each bus — Bus~1 is the reference ($\theta_1 = 0^{\circ}$), with angles of $-9.55^{\circ}$ and $-13.37^{\circ}$ at Buses~2 and~3;
    (b) LOPF generator dispatch for Scenario~A (unconstrained, blue) and Scenario~B (200\,MW thermal limit, red);
    (c) transmission line power flows for both scenarios, with the 200\,MW thermal limit shown as a dashed line — Line~1-3 is at 116.7\% loading in Scenario~A and exactly at the limit in Scenario~B.}
    \label{fig:dc_pf_result}
\end{figure}

% ──────────────────────────────────────────────────────────────────────────────
\section{AC Power Flow Implementation}
% ──────────────────────────────────────────────────────────────────────────────

\subsection{PowerModels.jl Data Format}

The AC power flow is implemented using PowerModels.jl, which requires a specific data dictionary format. Building this format from scratch revealed several non-obvious requirements:

\begin{enumerate}
    \item \textbf{Per-unit conversion:} All quantities must be in p.u.\ ($S_\text{base} = 100\,\text{MVA}$). Generator output, load power, and line impedances must all be divided by their respective base values. Failing to convert leads to angles of hundreds or thousands of radians in the solver output.

    \item \textbf{Mandatory shunt admittance fields:} Each branch requires explicit \texttt{g\_fr}, \texttt{b\_fr}, \texttt{g\_to}, \texttt{b\_to} fields (shunt conductance and susceptance at the from- and to-ends of the line). These are not auto-computed; for lossless lines, \texttt{g\_fr = g\_to = 0} and \texttt{b\_fr = b\_to = $b_c/2$}, where $b_c$ is the line charging susceptance. This requirement was discovered by reading the PowerModels.jl source (\texttt{matpower.jl:454}).

    \item \textbf{Angle bounds in radians:} The \texttt{angmin}/\texttt{angmax} fields must be in radians ($\pm\pi/3 \approx \pm60^{\circ}$). Supplying values in degrees makes the bounds effectively infinite.

    \item \textbf{Branch flow extraction:} The \texttt{build\_pf} formulation uses JuMP affine expressions (not variables) for branch flows, so they do not appear in the solution dictionary. Branch flows must be computed post-solution via the helper \texttt{calc\_branch\_flow\_ac} from PowerModels.
\end{enumerate}

\subsection{Implementation}

The key function converting the network to PowerModels format is shown below (condensed for clarity):

\begin{lstlisting}[language=Julia, caption={PowerModels.jl data format construction (key fields)}, label={lst:ac_format}]
function network_to_powermodels(buses, lines, generators, loads,
                                 v_nom=380.0; baseMVA=100.0)
    data = Dict{String,Any}(
        "baseMVA"  => baseMVA,
        "per_unit" => true,
        "shunt"    => Dict{String,Any}(),   # required empty dicts
        "dcline"   => Dict{String,Any}(),
        "storage"  => Dict{String,Any}(),
        "switch"   => Dict{String,Any}(), ...
    )
    # Branch: convert to per-unit + add shunt admittance fields
    Z_base = v_nom^2 / baseMVA          # base impedance [Ohm]
    br_r = r / Z_base;  br_x = x / Z_base
    br_b = 1.0 / sqrt(br_r^2 + br_x^2) * 0.1  # ~10% charging
    branch_data = Dict(
        "br_r"  => br_r,  "br_x"  => br_x,
        "g_fr"  => 0.0,   "b_fr"  => br_b/2.0,   # required!
        "g_to"  => 0.0,   "b_to"  => br_b/2.0,   # required!
        "angmin" => -pi/3, "angmax" => pi/3, ...)
    # Generator: per-unit active power
    gen_data = Dict("pg" => P_MW / baseMVA,
                    "pmax" => P_max / baseMVA, ...)
end
\end{lstlisting}

\subsection{Branch Flow Extraction}

After solving, branch flows are recovered in two steps: the solution is merged back into the data dictionary, and then the branch flows are computed from the recovered voltages:

\begin{lstlisting}[language=Julia, caption={Extracting branch flows from PowerModels solution}, label={lst:branch_flows}]
PowerModels.update_data!(pm_data, result["solution"])
branch_flows = PowerModels.calc_branch_flow_ac(pm_data)
for (i, branch) in sort(branch_flows["branch"])
    pf_MW = branch["pf"] * baseMVA     # convert from p.u. to MW
    pt_MW = branch["pt"] * baseMVA
end
\end{lstlisting}

\subsection{Validation}

The AC power flow implementation is validated on the same 3-bus test network defined in Section~3.1, with the generator dispatch fixed at $G_1 = 500\,\text{MW}$ (slack), loads $L_2 = 300\,\text{MW}$ and $L_3 = 200\,\text{MW}$, and line parameters $r = 0.01\,\text{p.u.}$, $x = 0.1\,\text{p.u.}$ ($S_\text{base} = 100\,\text{MVA}$, $V_\text{base} = 380\,\text{kV}$).

Table~\ref{tab:ac_validation} compares the Julia (PowerModels.jl + Ipopt) and PyPSA (\texttt{network.pf()}) results.

\begin{table}[H]
\centering
\caption{AC power flow validation: Julia vs PyPSA (3-bus network)}
\label{tab:ac_validation}
\begin{tabular}{lrrr}
\toprule
\textbf{Quantity} & \textbf{Julia} & \textbf{PyPSA} & \textbf{Max $|\Delta|$} \\
\midrule
$|V_1|$ (p.u.)    & 1.000000  & 1.000000  & 0 \\
$|V_2|$ (p.u.)    & 0.999982  & 0.999982  & $4.8 \times 10^{-7}$ \\
$|V_3|$ (p.u.)    & 0.999984  & 0.999984  & $4.8 \times 10^{-7}$ \\
$\theta_1$ (rad)  & 0.000000  & 0.000000  & 0 \\
$\theta_2$ (rad)  & $-$0.000185 & $-$0.000185 & $4.1 \times 10^{-7}$ \\
$\theta_3$ (rad)  & $-$0.000162 & $-$0.000162 & $4.1 \times 10^{-7}$ \\
$p_{12}$ (MW)     & 266.67    & 266.67    & $3.5 \times 10^{-3}$ \\
$p_{13}$ (MW)     & 233.34    & 233.34    & $3.5 \times 10^{-3}$ \\
$p_{23}$ (MW)     & $-$33.33  & $-$33.33  & $3.5 \times 10^{-3}$ \\
\bottomrule
\end{tabular}
\end{table}

All quantities agree with PyPSA to better than $5 \times 10^{-3}$ absolute error. The small residual in line active power ($3.5 \times 10^{-3}\,\text{MW}$) reflects the difference between Ipopt's interior-point convergence tolerance ($10^{-8}\,\text{p.u.}$) and PyPSA's Newton--Raphson tolerance ($10^{-6}\,\text{p.u.}$), not a modelling error. Voltage magnitudes and angles agree to $< 5 \times 10^{-7}$, confirming that the per-unit conversion, admittance matrix construction, and PowerModels.jl data format are all correct.

% ──────────────────────────────────────────────────────────────────────────────
\section{Linearized AC Power Flow Implementation}
% ──────────────────────────────────────────────────────────────────────────────

\subsection{Implementation}

The Linearized AC Power Flow is implemented in \texttt{julia/solvers/linear\_ac\_pf.jl} as a direct linear solve of Equation~\eqref{eq:lacpf}. The key difference from the DC power flow is the construction of both the susceptance Laplacian $\mathbf{B}'$ and the conductance Laplacian $\mathbf{G}'$, and the assembly of the $2(n-1) \times 2(n-1)$ coupled system:

\begin{lstlisting}[language=Julia, caption={Linearized AC PF: admittance matrices and coupled solve}, label={lst:lacpf_core}]
z_base = v_nom^2 / baseMVA          # 380^2 / 100 = 1444 Ohm
for (from, to, r, x) in lines
    r_pu = r / z_base;  x_pu = x / z_base
    denom = r_pu^2 + x_pu^2
    g_ij = r_pu / denom;  b_ij = x_pu / denom
    # Susceptance Laplacian B'
    B[from,from] += b_ij;  B[to,to] += b_ij
    B[from,to]   -= b_ij;  B[to,from] -= b_ij
    # Conductance Laplacian G'
    G[from,from] += g_ij;  G[to,to] += g_ij
    G[from,to]   -= g_ij;  G[to,from] -= g_ij
end
# 2(n-1) x 2(n-1) coupled system
A = [ B_r  G_r ; -G_r  -B_r ]
x_sol = A \ [P_r; Q_r]       # single linear solve
dtheta, dv = x_sol[1:nm], x_sol[nm+1:end]
V_mag = 1.0 .+ dv;  theta = dtheta
\end{lstlisting}

The implementation requires no iterative solver --- a single backslash call (dispatching to LAPACK) solves the coupled active/reactive system. The computational cost is $O((2n)^3)$ for dense matrices, but the sparsity of real networks means sparse LU factorisation is used for large $n$.

\subsection{Validation}

Table~\ref{tab:lacpf_validation} compares the LACPF solution against the full AC power flow reference (PyPSA Newton--Raphson) and the DC power flow, on the standard 3-bus network.

\begin{table}[H]
\centering
\caption{Linearized AC PF validation: LACPF vs full AC PF reference (3-bus network)}
\label{tab:lacpf_validation}
\begin{tabular}{lrrrr}
\toprule
\textbf{Quantity} & \textbf{LACPF} & \textbf{Full AC (ref)} & \textbf{DC PF} & \textbf{$|\Delta|$ LACPF} \\
\midrule
$\theta_2$ (rad)    & $-$0.000185 & $-$0.000185 & $-$0.000185 & $3.3 \times 10^{-7}$ \\
$\theta_3$ (rad)    & $-$0.000162 & $-$0.000162 & $-$0.000162 & $4.1 \times 10^{-7}$ \\
$|V_2|$ (p.u.)      &  1.000018   &  0.999982   &  1.000000   & $3.7 \times 10^{-5}$ \\
$|V_3|$ (p.u.)      &  1.000016   &  0.999984   &  1.000000   & $3.2 \times 10^{-5}$ \\
$p_{12}$ (MW)       &  261.39     &  266.67     &  266.67     & 5.28\,MW \\
$q_{12}$ (MVAr)     & $-$52.81    &   0.048     & N/A         & 52.85\,MVAr \\
\bottomrule
\end{tabular}
\end{table}

\textbf{Voltage angles and magnitudes} agree with the full AC reference to $< 4 \times 10^{-5}$\,p.u.\ and $< 5 \times 10^{-7}$\,rad, respectively --- a substantial improvement over the DC approximation, which cannot estimate voltage magnitudes at all.

\textbf{Active power flows} show a 5.28\,MW discrepancy ($\approx 2\%$) compared to the full AC solution. This difference arises because the LACPF includes the resistive correction through $\mathbf{G}'$: in a lossless DC network the flow is 266.67\,MW, while the LACPF redistributes $\approx 5\,\text{MW}$ into resistive losses.

\textbf{Reactive power flows} exhibit a large error ($52.85\,\text{MVAr}$). This is expected: the flat-start linearisation assumes $Q^{(0)} = 0$, but the actual reactive power depends on second-order voltage effects that are entirely discarded at first order. For applications requiring accurate reactive power, the LACPF serves as a warm start for Newton--Raphson, not a replacement.

% ──────────────────────────────────────────────────────────────────────────────
\section{LOPF Implementation}
% ──────────────────────────────────────────────────────────────────────────────

\subsection{JuMP Model Structure}

The LOPF is implemented in \texttt{julia/lopf.jl} using JuMP.jl. The model formulation maps directly to Equations~\eqref{eq:lopf_obj}--\eqref{eq:lopf_gen}:

\begin{lstlisting}[language=Julia, caption={LOPF JuMP model (core formulation)}, label={lst:lopf_jump}]
function solve_lopf(buses, lines, generators, loads, line_names;
                    line_capacity=Inf, baseMVA=100.0)
    n_buses  = length(buses)
    gen_buses = sort(collect(keys(generators)))

    # Susceptance matrix [MW/rad]
    B = zeros(n_buses, n_buses)
    sus = Float64[]
    for (from, to, r, x) in lines
        b = baseMVA / x          # x in per-unit
        push!(sus, b)
        B[from,from]+=b; B[to,to]  +=b
        B[from,to]  -=b; B[to,from]-=b
    end

    model = Model(HiGHS.Optimizer);  set_silent(model)

    @variable(model, theta[1:n_buses])
    @variable(model, P_gen[bus in gen_buses],
              lower_bound=0.0,
              upper_bound=generators[bus][1])    # capacity

    @constraint(model, theta[gen_buses[1]] == 0.0)   # slack
    for k in 1:n_buses
        P_inj = (k in gen_buses) ? P_gen[k] : 0.0
        @constraint(model,
            sum(B[k,m]*theta[m] for m in 1:n_buses) == P_inj - P_load[k])
    end
    if isfinite(line_capacity)
        for (i,(from,to,_,_)) in enumerate(lines)
            f = sus[i] * (theta[from] - theta[to])
            @constraint(model, -line_capacity <= f <= line_capacity)
        end
    end
    @objective(model, Min,
        sum(generators[bus][2] * P_gen[bus] for bus in gen_buses))
    optimize!(model)
end
\end{lstlisting}

\subsection{Validation: 3-Bus LOPF}

Using the test network defined in Section~3.2, the LOPF was solved under two scenarios:

\begin{itemize}
    \item \textbf{Scenario~A} (unconstrained): No line capacity limits ($F_{ij}^\text{max} = \infty$).
    \item \textbf{Scenario~B} (congested): Line thermal limit $F_{ij}^\text{max} = 200\,\text{MW}$ on all lines.
\end{itemize}

\begin{table}[H]
\centering
\caption{LOPF validation: 3-bus network, two constraint scenarios (Julia vs PyPSA)}
\label{tab:lopf_validation}
\begin{tabular}{lccccl}
\toprule
\textbf{Scenario} & \textbf{$P_1$ (MW)} & \textbf{$P_2$ (MW)} & \textbf{Max flow (MW)} & \textbf{Cost (\euro/h)} & \textbf{Matches?} \\
\midrule
A: No limits        & 400 & 100 & 233 & 13{,}000 & Yes ($\Delta = 0$) \\
B: 200\,MW limit    & 300 & 200 & 200 & 16{,}000 & Yes ($\Delta = 0$) \\
\bottomrule
\end{tabular}
\end{table}

\textbf{Scenario~A} (unconstrained): The optimizer dispatches the cheap generator G1 ($c_1 = 20\,\text{\euro/MWh}$) as much as possible and uses G2 ($c_2 = 50\,\text{\euro/MWh}$) only to cover the residual:
\begin{equation}
    P_1 = 400\,\text{MW},\quad P_2 = 100\,\text{MW}, \quad
    \text{Cost} = 400 \times 20 + 100 \times 50 = 13{,}000\,\text{\euro/h}.
\end{equation}
The maximum line flow is $p_{13} = 233\,\text{MW}$ (line 1--3), which exceeds the 200\,MW thermal limit.

\textbf{Scenario~B} (200\,MW limit): The binding constraint on line 1--3 forces a different dispatch. G1 is reduced and G2 must compensate:
\begin{equation}
    P_1 = 300\,\text{MW},\quad P_2 = 200\,\text{MW}, \quad
    \text{Cost} = 300 \times 20 + 200 \times 50 = 16{,}000\,\text{\euro/h} \quad (+23.1\%).
\end{equation}
The 23.1\% cost increase directly quantifies the economic cost of network congestion. Both results match PyPSA to numerical precision ($\Delta = 0\,\text{MW}$, $\Delta = 0\,\text{\euro/h}$).

The dispatch and line flow results for both scenarios are shown in panels (b) and (c) of Figure~\ref{fig:dc_pf_result}.

% ──────────────────────────────────────────────────────────────────────────────
\section{Validation on the IEEE 14-Bus Standard Test Case}
\label{sec:ieee14}
% ──────────────────────────────────────────────────────────────────────────────

The three-bus network used in Sections~3.2--3.4 is intentionally simple and designed to allow analytical verification. To demonstrate that the implementations generalise to a recognised industry benchmark, all three modules were additionally validated on the \textbf{IEEE 14-bus test case} \cite{zimmerman2011matpower} --- a standard network used throughout the power systems literature for algorithm verification and comparison.

The IEEE 14-bus network comprises 14 buses, 20 branches (including 4 transformers with off-nominal tap ratios), 5 generators, and 11 load buses with a total active demand of 259\,MW ($S_\text{base} = 100\,\text{MVA}$). The MATPOWER case file \texttt{case14.m} was used as the common data source for both Julia and Python, ensuring that both implementations operate on identical network data.

\subsection{AC Power Flow on IEEE 14-Bus}

Table~\ref{tab:ieee14_ac} compares the Julia (PowerModels.jl + Ipopt) AC power flow solution against the reference solution stored in \texttt{case14.m}, which represents the MATPOWER-solved steady-state. PowerModels.jl reads and models the transformers with their correct tap ratios; the Julia implementation therefore produces a full transformer-aware solution.

\begin{table}[H]
\centering
\caption{IEEE 14-bus AC power flow: Julia vs MATPOWER reference}
\label{tab:ieee14_ac}
\begin{tabular}{rrrrrr}
\toprule
\textbf{Bus} & \textbf{$|V|$ Julia (p.u.)} & \textbf{$|V|$ Ref} & \textbf{$|\Delta V|$} & \textbf{$\theta$ Julia (°)} & \textbf{$|\Delta\theta|$ (°)} \\
\midrule
 1 & 1.0600 & 1.0600 & $0$          &   0.0000 & $0$        \\
 2 & 1.0450 & 1.0450 & $0$          &  $-$4.9826 & $2.6\times10^{-3}$ \\
 3 & 1.0100 & 1.0100 & $0$          & $-$12.7251 & $5.1\times10^{-3}$ \\
 4 & 1.0177 & 1.0190 & $1.3\times10^{-3}$ & $-$10.3129 & $1.7\times10^{-2}$ \\
 6 & 1.0700 & 1.0700 & $0$          & $-$14.2209 & $9.5\times10^{-4}$ \\
 9 & 1.0559 & 1.0560 & $6.8\times10^{-5}$ & $-$14.9385 & $1.5\times10^{-3}$ \\
14 & 1.0355 & 1.0360 & $4.7\times10^{-4}$ & $-$16.0336 & $6.4\times10^{-3}$ \\
\midrule
\multicolumn{3}{l}{\textbf{Max over all 14 buses}} & $\mathbf{1.33\times10^{-3}}$ & & $\mathbf{1.71\times10^{-2}}$ \\
\bottomrule
\end{tabular}
\end{table}

The maximum voltage magnitude error is $1.33\times10^{-3}$\,p.u.\ and the maximum angle error is $0.017^{\circ}$, both well within engineering accuracy ($<0.5\%$). The small residual reflects the difference in convergence tolerances between Ipopt and the reference solver, not a modelling error. The result confirms that the PowerModels.jl-based AC power flow implementation handles transformer tap ratios correctly on a realistic multi-machine network.

\subsection{DC Power Flow on IEEE 14-Bus}

Table~\ref{tab:ieee14_dc} reports the DC power flow voltage angles from Julia and PyPSA, both compared against the full AC reference.

\begin{table}[H]
\centering
\caption{IEEE 14-bus DC power flow: Julia and PyPSA vs full-AC reference (MATPOWER)}
\label{tab:ieee14_dc}
\begin{tabular}{rrrr}
\toprule
\textbf{Bus} & \textbf{$\theta$ Julia (°)} & \textbf{$\theta$ PyPSA (°)} & \textbf{$|\Delta\theta|$ vs AC ref (°)} \\
\midrule
 1 &   0.000 &   0.000 & 0.000 \\
 2 &  $-$5.013 &  $-$5.013 & 0.033 \\
 5 &  $-$9.089 &  $-$9.089 & 0.309 \\
 6 & $-$15.165 & $-$15.165 & 0.945 \\
 9 & $-$15.893 & $-$15.893 & 0.953 \\
14 & $-$17.429 & $-$17.429 & 1.389 \\
\midrule
\multicolumn{2}{l}{\textbf{Max (all 14 buses)}} & \textbf{identical} & \textbf{1.389°} \\
\bottomrule
\end{tabular}
\end{table}

Julia and PyPSA produce \emph{identical} DC power flow angles across all 14 buses. The deviation from the full AC reference (max $1.39^{\circ}$) is larger than for pure transmission networks because IEEE 14-bus contains four transformers with off-nominal tap ratios ($0.932$--$0.978$), which the DC linearisation ignores. This is a known limitation of the DC approximation, not a software defect.

\subsection{LOPF on IEEE 14-Bus}

Both implementations minimise total generation cost over the five generators at buses 1, 2, 3, 6, and 8. The two cheapest generators (buses 1 and 2, marginal cost $20\,\text{\euro/MWh}$) are preferred over the more expensive generators (buses 3, 6, 8 at $40\,\text{\euro/MWh}$).

\begin{table}[H]
\centering
\caption{IEEE 14-bus LOPF: Julia (JuMP+HiGHS) vs PyPSA (linopy+HiGHS)}
\label{tab:ieee14_lopf}
\begin{tabular}{lrrrr}
\toprule
\textbf{Metric} & \textbf{Julia} & \textbf{PyPSA} & \textbf{Match?} \\
\midrule
Total load (MW)     & 259.0 & 259.0 & Yes \\
Total cost (\euro/h) & 5{,}180 & 5{,}180 & Yes \\
Generators 3, 6, 8  & 0\,MW each & 0\,MW each & Yes \\
\bottomrule
\end{tabular}
\end{table}

Both solvers find the same optimal objective value of $5{,}180\,\text{\euro/h}$ and correctly set the three expensive generators to zero. The LP is degenerate (generators 1 and 2 share the same marginal cost of $20\,\text{\euro/MWh}$), so the individual dispatch at buses 1 and 2 may differ between solvers while the total cost and expensive-generator dispatch remain identical --- a well-known property of degenerate LP optima.

% ──────────────────────────────────────────────────────────────────────────────
\section{Multi-Period LOPF with Storage and Wind}
\label{sec:multiperiod}
% ──────────────────────────────────────────────────────────────────────────────

The single-period LOPF of Section~3.4 dispatches generators for a single snapshot in time. Real-world energy system planning requires optimising dispatch across \emph{multiple} time periods simultaneously, so that cheap energy can be stored when it is abundant and released when demand is high --- the classical \emph{time-arbitrage} problem. This section extends the Julia LOPF to a 24-hour multi-period formulation with a \textbf{StorageUnit} and a \textbf{variable wind generator}, and validates it against PyPSA's built-in multi-period optimisation.

\subsection{Mathematical Extension}

The multi-period LOPF adds a time index $t \in \{1, \ldots, T\}$ (with $T = 24$\,h) to the single-period LP. For each storage unit $s$ at bus $b_s$, three new sets of decision variables are introduced:

\begin{itemize}
    \item $p_{s,t}^{ch}$ --- charging power [MW], constrained to $[0, P_s^{nom}]$
    \item $p_{s,t}^{dis}$ --- discharging power [MW], constrained to $[0, P_s^{nom}]$
    \item $e_{s,t}$ --- state of charge (SOC) [MWh], constrained to $[0, E_s^{nom}]$
\end{itemize}

The SOC evolves according to the energy balance equation:
\begin{equation}
    e_{s,t} = e_{s,t-1} + \eta_s^{ch}\,p_{s,t}^{ch} - \frac{p_{s,t}^{dis}}{\eta_s^{dis}},
    \quad \forall t \in \{1,\ldots,T\},
    \label{eq:soc_dynamics}
\end{equation}
where $\eta_s^{ch}$ and $\eta_s^{dis}$ are the charge and discharge efficiencies. A \emph{cyclicity} constraint enforces that the SOC at the end of the planning horizon equals the initial SOC:
\begin{equation}
    e_{s,T} = e_{s,0}.
    \label{eq:soc_cyclic}
\end{equation}
This prevents the optimiser from artificially depleting storage at the end of the horizon. The wind generator contributes a fixed, time-varying active power injection $P_{w,t}^{wind} = p_{w,t}^{max} \cdot P_w^{nom}$, where $p_{w,t}^{max} \in [0,1]$ is the capacity factor at time $t$. Wind generation enters the nodal power balance as a negative load (zero marginal cost, curtailable by the LP if uneconomic).

The full objective and power balance constraints are:
\begin{align}
    \min \quad & \sum_{t=1}^{T} \sum_{i \in \mathcal{G}} c_i\,P_{i,t}^{gen}
    \label{eq:mp_obj} \\
    \text{s.t.} \quad
    & \sum_{j} B_{kj}\,\theta_{j,t} = P_{k,t}^{gen} + P_{k,t}^{wind}
      + p_{s(k),t}^{dis} - p_{s(k),t}^{ch} - P_{k,t}^{load},
      \quad \forall k,t \label{eq:mp_balance}
\end{align}
together with generator limits, line limits, and Equations~\eqref{eq:soc_dynamics}--\eqref{eq:soc_cyclic}.

\subsection{Test Scenario and Results}

The 3-bus network from Section~3.1 is extended with a wind generator at Bus~3 and a storage unit at Bus~2. Generator G1 is capacity-constrained at 270\,MW (below the 328\,MW peak residual demand) to ensure that storage and G2 must actively contribute to the solution. Table~\ref{tab:mp_params} summarises the component parameters.

\begin{table}[H]
\centering
\caption{Multi-period LOPF: component parameters}
\label{tab:mp_params}
\begin{tabular}{llll}
\toprule
\textbf{Component} & \textbf{Bus} & \textbf{Capacity} & \textbf{Cost / Parameters} \\
\midrule
Generator G1   & 1 & 270\,MW        & 20\,\euro/MWh \\
Generator G2   & 2 & 100\,MW        & 50\,\euro/MWh \\
Storage        & 2 & 100\,MW / 300\,MWh & $\eta^{ch} = \eta^{dis} = 0.95$ \\
Wind           & 3 & 150\,MW nom.   & 0\,\euro/MWh, $p^{max}(t)$ variable \\
Load (base)    & 2 & 250\,MW        & daily profile $\times$ [0.54--1.00] \\
Load (base)    & 3 & 175\,MW        & same profile, peak total = 425\,MW \\
\bottomrule
\end{tabular}
\end{table}

Both Julia (JuMP + HiGHS) and PyPSA (linopy + HiGHS) find the same optimal cost:

\begin{equation}
    \text{Total cost (24\,h)} = 126{,}790.79\,\text{\euro},
    \qquad
    \text{Average} = 15.65\,\text{\euro/MWh}.
    \label{eq:mp_cost}
\end{equation}

Table~\ref{tab:mp_dispatch} shows the hourly dispatch for selected hours, illustrating the economic logic of storage operation.

\begin{table}[H]
\centering
\caption{Multi-period LOPF: selected hourly dispatch (Julia solution)}
\label{tab:mp_dispatch}
\begin{tabular}{rrrrrrl}
\toprule
\textbf{Hour} & \textbf{Load} & \textbf{Wind} & \textbf{G1} & \textbf{G2} & \textbf{Storage} & \textbf{Action} \\
& (MW) & (MW) & (MW) & (MW) & (MW) & \\
\midrule
 3 &  229.5 & 124.5 & 150.4 &  0.0 & $-$45.4 & charging \\
 4 &  233.8 & 117.0 & 216.8 &  0.0 & $-$100.0 & charging (max) \\
 9 &  382.5 &  67.5 & 270.0 &  0.0 &  45.0 & discharging \\
10 &  391.0 &  63.0 & 270.0 & 38.8 &  19.2 & discharging + G2 \\
18 &  425.0 &  97.5 & 270.0 & 57.5 &   0.0 & G2 only (storage depleted) \\
\bottomrule
\end{tabular}
\end{table}

The dispatch exhibits classical time-arbitrage behaviour: the storage charges during low-demand night hours using cheap G1 energy ($\approx 22.2\,\text{\euro/MWh}$ round-trip cost including $\eta^2 = 0.9025$ efficiency) and discharges during peak hours to displace expensive G2 ($50\,\text{\euro/MWh}$). Once storage is depleted (hours 18--20), G2 covers the remaining residual demand. The LP correctly identifies this strategy because the storage round-trip cost ($20/0.9025 \approx 22.2\,\text{\euro/MWh}$) is well below G2's marginal cost ($50\,\text{\euro/MWh}$).

\textbf{Validation:} the Julia and PyPSA solutions agree to the cent on total cost. Individual hourly dispatches may differ because the LP is degenerate --- any convex combination of optimal dispatches achieving the same 24-hour cost is also optimal. The identical objective value confirms mathematical correctness.

\subsection{Multi-Period LOPF Performance}

Table~\ref{tab:mp_benchmark} reports the multi-period LOPF solve times as a function of the planning horizon $T$, keeping the 3-bus network fixed. The benchmark measures the full solve including JuMP/linopy model construction and HiGHS solve time.

\begin{table}[H]
\centering
\caption{Multi-period LOPF benchmark: Julia (JuMP+HiGHS) vs PyPSA (linopy+HiGHS), 3-bus network, minimum time}
\label{tab:mp_benchmark}
\begin{tabular}{rrrr}
\toprule
\textbf{T (hours)} & \textbf{Julia (ms)} & \textbf{PyPSA (ms)} & \textbf{Speedup} \\
\midrule
  6 &  3.79 & 1{,}131 & 298$\times$ \\
 12 &  5.37 & 1{,}151 & 214$\times$ \\
 24 &  8.74 & 1{,}141 & 131$\times$ \\
 48 & 12.82 & 1{,}136 &  89$\times$ \\
\bottomrule
\end{tabular}
\end{table}

Two observations stand out. First, the PyPSA solve time is nearly \emph{constant} across all horizons ($\approx 1.14\,\text{s}$), independent of whether the LP spans 6 or 48 periods. This confirms that PyPSA's fixed overhead --- pandas network state management, linopy model assembly, and Python object allocation --- completely dominates the actual LP solve time at this network size; the HiGHS solver itself completes in well under a millisecond, and the remaining $\sim 1.14\,\text{s}$ is model-building overhead.

Second, Julia's solve time grows approximately \emph{linearly} with $T$: from $3.8\,\text{ms}$ at $T = 6$ to $12.8\,\text{ms}$ at $T = 48$. This near-linear scaling reflects JuMP's compiled sparse constraint assembly, where adding $T$ periods adds $T$ independent constraint sets with no global restructuring.

The consequence is that the speedup \emph{decreases} with $T$ (from $298\times$ to $89\times$). This does not indicate diminishing returns for Julia --- PyPSA remains $89\times$ slower even at $T = 48$ --- but rather that PyPSA's cost is a large fixed model-assembly term that Julia's growing per-period cost only slowly erodes.

% ──────────────────────────────────────────────────────────────────────────────
\section{Comprehensive Component Validation}
\label{sec:full_validation}
% ──────────────────────────────────────────────────────────────────────────────

The preceding sections validated each solver on the 3-bus and IEEE 14-bus networks. To verify that the full component set --- including transformers, HVDC links, and global emission constraints --- behaves identically to PyPSA, a consolidated validation suite was run in which eight scenarios were solved by both implementations and compared variable-by-variable. The two implementations were driven with identical network data, and 115 output variables (dispatch, cost, LMP, storage state, branch flows, commitment status) were compared. Table~\ref{tab:full_validation} summarises the outcome.

\begin{table}[H]
\centering
\caption{Comprehensive validation: Julia vs PyPSA across all components (115 variables, 8 scenarios)}
\label{tab:full_validation}
\begin{tabular}{llc}
\toprule
\textbf{Scenario} & \textbf{Variables compared} & \textbf{Agreement} \\
\midrule
LOPF uncongested          & $P_g$, cost, LMP            & exact ($\Delta = 0$) \\
LOPF congested            & $P_g$, cost, LMP            & exact ($\Delta = 0$) \\
LOPF + Transformer        & $P_g$, cost, LMP            & exact ($\Delta = 0$) \\
LOPF + HVDC Link          & $P_g$, $P_\text{HVDC}$, cost, LMP & exact ($\Delta = 0$) \\
LOPF + CO$_2$ constraint  & $P_g$, cost, LMP            & exact ($\Delta = 0$) \\
Multi-period + StorageUnit & $P_g(t)$, SoC$(t)$, LMP$(t)$, cost & exact ($\Delta = 0$) \\
DC Power Flow             & branch flows                & exact ($\Delta = 0$) \\
DC Power Flow             & voltage angles              & scale factor (convention) \\
LOPF + Transformer        & transformer branch flow     & scale factor (convention) \\
Unit Commitment           & dispatch, cost, commitment  & Julia optimal (PyPSA artifact) \\
\midrule
\multicolumn{2}{l}{\textbf{Economically meaningful quantities (dispatch, cost, LMP)}} & \textbf{100\% exact} \\
\bottomrule
\end{tabular}
\end{table}

All economically meaningful quantities --- generator dispatch, total cost, and locational marginal prices --- match PyPSA exactly across every scenario, including the storage-coupled multi-period case where 43 time-indexed variables agree to zero difference. Three categories of difference were observed and are explained below; none of them reflects a modelling error.

\textbf{DC power flow voltage angles} differ by a constant factor of $V_\text{nom}^2 / S_\text{base} = 380^2/100 = 1444$. PyPSA expresses the susceptance matrix in per-unit on the nominal voltage base, producing angles on the order of $10^{-4}$\,rad, whereas the present implementation follows the MATPOWER convention $b = S_\text{base}/x$ in MW/rad, producing angles on the order of $10^{-1}$\,rad. The two are related by an exact scalar and yield identical branch flows (which were verified to match to $\Delta = 0$); the choice of angle scale is a presentation convention, not a physical difference.

\textbf{The transformer branch flow} differs because PyPSA refers the transformer reactance to the transformer's own rating $s_\text{nom}$, whereas the present implementation refers it to the system base. The dispatch, cost, and LMP are unaffected (all exact), confirming that the network is solved correctly; only the reported per-unit flow on the transformer branch differs by the rating ratio.

\textbf{The Unit Commitment dispatch} differs in one scenario, and here the Julia solver reaches a \emph{lower} total cost than PyPSA ($34{,}100$ versus $36{,}800$). The Julia solution is the verifiable optimum: the cheap base generator (400~MW at \euro20/MWh) alone covers the peak demand of 300~MW, so committing the peaker is never economical. PyPSA nevertheless commits the peaker at $90$~MW in the first period. A controlled re-run (\texttt{python/recheck\_uc.py}) isolates the cause to the interaction of \texttt{min\_up\_time}~$=2$ with \texttt{p\_min\_pu}~$=0.3$: setting either \texttt{min\_up\_time}~$=1$ or \texttt{p\_min\_pu}~$=0$ makes PyPSA return the same $34{,}100$ optimum, whereas \texttt{initial\_status} and the slack-bus setting have no effect. The PyPSA schedule moreover keeps the peaker on for a single period, which violates its own \texttt{min\_up\_time}~$=2$ --- a boundary artifact of PyPSA's min-up-time formulation at the start of the horizon, not a discrepancy in the Julia model. Finally, the locational marginal prices differ because a mixed-integer program has no dual solution: PyPSA reports zero prices for this committable-only case, while the Julia implementation recovers meaningful prices by fixing the binaries and re-solving the resulting LP.

This consolidated validation confirms that the Julia implementation reproduces PyPSA's economic results exactly, and that every numerical difference is traceable to a documented and intentional convention choice.

% ──────────────────────────────────────────────────────────────────────────────
\section{Performance Benchmark Results}
% ──────────────────────────────────────────────────────────────────────────────

\subsection{Benchmark Methodology}
\label{sec:bench_methodology}

All timings reported in this section follow a single, fair protocol designed to isolate the language-level contribution to performance. For every scenario the network is constructed \emph{once}; only the per-solve operation --- model assembly plus solve --- is timed, on both the Julia and the PyPSA side, so that one-time network-construction cost is excluded from both. Each timed call is preceded by a warm-up call that excludes Julia's just-in-time compilation. The Julia times use \texttt{BenchmarkTools.jl} (auto-tuned repeated samples); the PyPSA times use repeated \texttt{perf\_counter} samples. To keep the comparison fair, the HiGHS solver and the BLAS linear-algebra backend are pinned to a \emph{single thread} on both sides. Throughout, the reported figure is the \textbf{minimum} over the samples, which is the standard noise-robust estimator of intrinsic compute time --- operating-system jitter and garbage-collection pauses can only add time, never subtract it. Synthetic networks are randomly generated with fixed seeds; in addition, standard published networks (IEEE and PEGASE) are benchmarked in Section~\ref{sec:bench_real}.

\subsection{DC Power Flow Results}

Table~\ref{tab:dc_benchmark} presents the DC power flow benchmark on synthetic networks. Julia exposes \texttt{pf(net)}; PyPSA uses \texttt{lpf()}. Both reduce to a sparse linear solve.

\begin{table}[H]
\centering
\caption{DC power flow benchmark on synthetic networks: Julia vs PyPSA (minimum time)}
\label{tab:dc_benchmark}
\begin{tabular}{rrrr}
\toprule
\textbf{Buses} & \textbf{Julia (ms)} & \textbf{PyPSA (ms)} & \textbf{Speedup} \\
\midrule
    3 &  0.002 & 123.61 & 61{,}803$\times$ \\
   10 &  0.006 & 118.13 & 21{,}088$\times$ \\
   50 &  0.057 & 123.93 &  2{,}174$\times$ \\
  100 &  0.171 & 163.64 &     957$\times$ \\
  300 &  2.135 & 360.83 &     169$\times$ \\
\bottomrule
\end{tabular}
\end{table}

The speedup is largest for small networks ($\sim$$6\times10^{4}$ at 3 buses) and falls steadily with size ($169\times$ at 300 buses). The very large small-network ratios are \emph{not} a measure of compute superiority: they reflect PyPSA's fixed per-call Python overhead ($\approx 120\,\text{ms}$) dwarfing a sub-millisecond Julia solve. The compute-bound figure --- recovered on large real networks in Section~\ref{sec:bench_real} --- is about $5\times$. This behavior is analyzed in Chapter~\ref{sec:analysis}.

\subsection{AC Power Flow Results}

Table~\ref{tab:ac_benchmark} presents the AC power flow benchmark results for networks of 3 to 100 buses. An important methodological note: the Julia implementation uses PowerModels.jl with the Ipopt interior-point solver, whereas PyPSA's \texttt{pf()} method applies a Newton--Raphson (NR) iteration directly to the AC power flow equations. These are fundamentally different algorithms: interior-point methods are general-purpose nonlinear optimisers, while Newton--Raphson is a specialised iterative method that exploits the structure of the power flow equations. Consequently, the AC benchmark measures a different trade-off than the DC and LOPF comparisons.

\begin{table}[H]
\centering
\caption{AC power flow benchmark: Julia (PowerModels.jl + Ipopt) vs PyPSA (Newton--Raphson), minimum time}
\label{tab:ac_benchmark}
\begin{tabular}{rrrr}
\toprule
\textbf{Buses} & \textbf{Julia (ms)} & \textbf{PyPSA (ms)} & \textbf{Speedup} \\
\midrule
   3 &   7.22 & 177.51 & 24.6$\times$ \\
  10 &  29.90 & 188.34 &  6.3$\times$ \\
  50 & 151.13 & 225.20 &  1.5$\times$ \\
 100 & 548.96 & 264.22 &  \textbf{0.5$\times$} \\
\bottomrule
\end{tabular}
\end{table}

AC power flow is the one method where Julia does \emph{not} hold a consistent advantage. The apparent speedup at small sizes is again dominated by PyPSA's fixed Python overhead ($\approx 180\,\text{ms}$); as that overhead is amortised, the comparison reverses, and at 100 buses Julia is about \textbf{twice as slow} as PyPSA ($0.5\times$). The reason is algorithmic, not linguistic: Julia delegates AC power flow to the Ipopt interior-point solver (via PowerModels.jl), whose per-iteration cost is higher than the specialised Newton--Raphson iteration that PyPSA's \texttt{pf()} applies directly to the mismatch equations. This is reported honestly: the model-assembly advantage that drives the LP results does not apply when the work is dominated by a third-party nonlinear solver. The implication --- that a native Newton--Raphson AC solver in Julia would be needed to recover an advantage --- is discussed in Chapter~\ref{sec:analysis_ac}.

\subsection{LOPF Results}

Table~\ref{tab:lopf_benchmark} presents the LOPF benchmark on synthetic networks. Both implementations use HiGHS as the LP solver backend; with the network built once and threads pinned, the measured difference is entirely in optimisation-model assembly, not in the solver itself.

\begin{table}[H]
\centering
\caption{LOPF benchmark on synthetic networks: Julia (JuMP+HiGHS) vs PyPSA (linopy+HiGHS), minimum time}
\label{tab:lopf_benchmark}
\begin{tabular}{rrrr}
\toprule
\textbf{Buses} & \textbf{Julia (ms)} & \textbf{PyPSA (ms)} & \textbf{Speedup} \\
\midrule
    3 &  2.28 &   645.69 & 283$\times$ \\
   10 &  2.64 &   669.13 & 253$\times$ \\
   50 &  4.62 &   844.997 & 183$\times$ \\
  100 &  6.39 & 1{,}043.92 & 163$\times$ \\
  300 & 17.08 & 2{,}022.00 & 118$\times$ \\
\bottomrule
\end{tabular}
\end{table}

The LOPF speedup is large and decreases gently with size (from $283\times$ at 3 buses to $118\times$ at 300 buses); it stabilises rather than collapsing, because PyPSA's cost is dominated by the Python/linopy model assembly that grows with the problem, whereas Julia assembles the same model in compiled code. The trend continues on large real networks below.

\subsection{Standard Real Networks (IEEE and PEGASE)}
\label{sec:bench_real}

To address the realism and scale of the synthetic study, the same two methods were benchmarked on standard published networks: IEEE~118 and~300, and the PEGASE~1354 and~2869 cases, which represent parts of the real European transmission grid \cite{josz2016ac}. The MATPOWER \texttt{.m} cases are loaded into a \texttt{PowerFlowJulia} network through the exported \texttt{add!} API and solved through \texttt{pf}/\texttt{optimize}, so the timing reflects the public interface rather than a hand-tuned kernel.

\begin{table}[H]
\centering
\caption{Benchmark on standard real networks (minimum time, ms): Julia vs PyPSA}
\label{tab:real_benchmark}
\begin{tabular}{lrrrrrr}
\toprule
 & \multicolumn{3}{c}{\textbf{DC Power Flow}} & \multicolumn{3}{c}{\textbf{LOPF}} \\
\cmidrule(lr){2-4}\cmidrule(lr){5-7}
\textbf{Network} & \textbf{Julia} & \textbf{PyPSA} & \textbf{$\times$} & \textbf{Julia} & \textbf{PyPSA} & \textbf{$\times$} \\
\midrule
IEEE 118     &   0.33 &   170.5 & 517 &   9.96 &  1{,}163 & 117 \\
IEEE 300     &   1.65 &   267.8 & 162 &  19.24 &  1{,}839 &  96 \\
PEGASE 1354  &  59.42 & 1{,}026 &  17 & 129.20 &  7{,}583 &  59 \\
PEGASE 2869  & 422.96 & 2{,}005 & 4.7 & 491.12 & 28{,}427 &  58 \\
\bottomrule
\end{tabular}
\end{table}

The flagship result is the 2{,}869-bus PEGASE LOPF: PyPSA takes $28.4\,\text{s}$, Julia $0.49\,\text{s}$, a $58\times$ speedup on a realistic European-scale grid. The LOPF speedup stays in the $58$--$117\times$ band across all four real networks, confirming that the model-assembly advantage persists at scale. The DC power flow speedup falls to $4.7\times$ at 2{,}869 buses --- the honest compute-bound figure once PyPSA's fixed overhead is amortised.

\subsection{Multi-Period LOPF and Unit Commitment}

The two time-coupled methods were also benchmarked against the planning horizon $T$ and, for unit commitment, against network size. To avoid duplication their results are reported in the dedicated sections: multi-period LOPF in Table~\ref{tab:mp_benchmark} (speedup $298\times \to 89\times$ as $T$ grows from 6 to 48, with PyPSA's time essentially constant at $\approx 1.14\,\text{s}$ --- a direct demonstration that the bottleneck is model assembly, not solving), and unit commitment, a mixed-integer program solved by the same HiGHS branch-and-bound on both sides, in Tables~\ref{tab:uc_benchmark_t} and~\ref{tab:uc_benchmark_scale} (speedup $17$--$79\times$).

\subsection{Benchmark Visualizations}

The benchmark figures (execution-time scaling on log--log axes, speedup curves, the real-network comparison, and the multi-period and unit-commitment panels) are produced by \texttt{make\_benchmark\_figures.py} from the result CSVs.

\begin{figure}[H]
    \centering
    \includegraphics[width=\textwidth]{plots/fig_scaling_dcpf_lopf.png}
    \caption{Execution-time scaling of (a)~DC power flow and (b)~LOPF on synthetic networks, Julia versus PyPSA, on log--log axes. Each point is the minimum solve time over repeated samples with threads pinned.}
    \label{fig:scaling_dcpf_lopf}
\end{figure}

\begin{figure}[H]
    \centering
    \includegraphics[width=0.7\textwidth]{plots/fig_speedup_size.png}
    \caption{Speedup of Julia over PyPSA (PyPSA time divided by Julia time) versus network size, for DC power flow and LOPF. The dashed line marks parity.}
    \label{fig:speedup_size}
\end{figure}

\begin{figure}[H]
    \centering
    \includegraphics[width=0.7\textwidth]{plots/fig_scaling_acpf.png}
    \caption{AC power flow execution time versus network size. Julia (Ipopt) is faster only at small sizes and becomes slower than PyPSA's Newton--Raphson at 100~buses --- the one method where the porting does not yield a speedup.}
    \label{fig:scaling_acpf}
\end{figure}

\begin{figure}[H]
    \centering
    \includegraphics[width=\textwidth]{plots/fig_scaling_mp_uc.png}
    \caption{Execution time versus planning horizon for (a)~multi-period LOPF and (b)~unit commitment. PyPSA's time is nearly constant, confirming that the bottleneck is model assembly rather than solving.}
    \label{fig:scaling_mp_uc}
\end{figure}

\begin{figure}[H]
    \centering
    \includegraphics[width=\textwidth]{plots/fig_realcases.png}
    \caption{Minimum solve time on standard IEEE and PEGASE networks for (a)~DC power flow and (b)~LOPF, Julia versus PyPSA. The PEGASE cases are part of the real European transmission grid.}
    \label{fig:realcases}
\end{figure}

% ──────────────────────────────────────────────────────────────────────────────
\section{Network Object Model and PyPSA-Compatible API}
\label{sec:network_model}
% ──────────────────────────────────────────────────────────────────────────────

The standalone solver scripts described in previous sections accept raw tuples and dictionaries as input, which is adequate for isolated benchmarks but impractical for building complete energy system models. To address this, a typed network object model was developed in \texttt{julia/src/}, providing a PyPSA-compatible API that allows the same workflow to be used in Julia and Python.

\subsection{Component Type Hierarchy}

All network components are represented as immutable Julia \texttt{struct}s, ensuring type stability and enabling compiler optimisations. Table~\ref{tab:components} lists the implemented components with their PyPSA equivalents.

\begin{table}[H]
\centering
\caption{PowerFlowJulia component types and PyPSA equivalents}
\label{tab:components}
\begin{tabular}{lll}
\toprule
\textbf{Julia Type} & \textbf{PyPSA Equivalent} & \textbf{Key Fields} \\
\midrule
\texttt{Bus}              & \texttt{network.buses}             & \texttt{v\_nom}, \texttt{slack}, \texttt{bus\_type} \\
\texttt{Line}             & \texttt{network.lines}             & \texttt{r}, \texttt{x}, \texttt{b}, \texttt{s\_nom} \\
\texttt{Transformer}      & \texttt{network.transformers}      & \texttt{x}, \texttt{tap\_ratio}, \texttt{phase\_shift} \\
\texttt{Generator}        & \texttt{network.generators}       & \texttt{p\_nom}, \texttt{marginal\_cost}, \texttt{committable} \\
\texttt{Load}             & \texttt{network.loads}            & \texttt{p\_set}, \texttt{q\_set} \\
\texttt{StorageUnit}      & \texttt{network.storage\_units}   & \texttt{p\_nom}, \texttt{e\_nom}, $\eta^{ch}$, $\eta^{dis}$ \\
\texttt{Store}            & \texttt{network.stores}           & \texttt{e\_nom}, \texttt{p\_nom}, \texttt{standing\_loss} \\
\texttt{Carrier}          & \texttt{network.carriers}         & \texttt{co2\_emissions} \\
\texttt{Link}             & \texttt{network.links}            & \texttt{bus0}, \texttt{bus1}, \texttt{efficiency} \\
\texttt{GlobalConstraint} & \texttt{network.global\_constraints} & \texttt{carrier\_weightings}, \texttt{constant} \\
\bottomrule
\end{tabular}
\end{table}

\subsection{Network Container and API}

The \texttt{Network} struct aggregates all component dictionaries (\texttt{Dict\{String,~T\}}) and exposes a unified \texttt{add!} interface that mirrors PyPSA's \texttt{network.add()} method. All ten component types --- \texttt{Bus}, \texttt{Line}, \texttt{Transformer}, \texttt{Generator}, \texttt{Load}, \texttt{StorageUnit}, \texttt{Store}, \texttt{Carrier}, \texttt{Link}, and \texttt{GlobalConstraint} --- are accessible through the same entry point, using the same keyword argument names as PyPSA (e.g., \texttt{bus0}/\texttt{bus1} for branch endpoints, \texttt{p\_nom}, \texttt{marginal\_cost}, \texttt{s\_nom}).

The high-level \texttt{pf} and \texttt{optimize} functions dispatch to the appropriate solver based on network content: if any generator has \texttt{committable=true}, \texttt{optimize} automatically selects Unit Commitment; if $T > 1$, it selects multi-period LOPF; otherwise, single-period LOPF is used. This design allows PyPSA users to migrate with minimal changes to their existing workflow.

% ──────────────────────────────────────────────────────────────────────────────
\section{Transformer Tap-Ratio Correction}
\label{sec:transformer}
% ──────────────────────────────────────────────────────────────────────────────

The IEEE 14-bus network contains four transformers with off-nominal tap ratios ($a \in \{0.932, 0.969, 0.978\}$). The standard DC power flow formulation ignores these ratios, treating transformers as ordinary lines. This introduces systematic angle errors at buses downstream of transformers.

\subsection{Corrected Transformer Model}

Following the PyPSA and MATPOWER simplified DC convention, the effective susceptance of a transformer with off-nominal tap ratio $a$ and reactance $x$ is:

\begin{equation}
    b_\text{eff} = \frac{S_\text{base}}{x \cdot a}.
    \label{eq:tap_correction}
\end{equation}

For phase-shifting transformers with shift angle $\varphi$ (degrees), the power balance right-hand side is modified by an equivalent injection:

\begin{equation}
    P_\text{shift}[k] += b_\text{eff} \cdot \varphi_\text{rad}, \qquad
    P_\text{shift}[m] -= b_\text{eff} \cdot \varphi_\text{rad},
    \label{eq:phase_shift}
\end{equation}

and the branch flow becomes $P_{km} = b_\text{eff}(\theta_k - \theta_m - \varphi_\text{rad})$.

\subsection{Validation on IEEE 14-Bus}

Table~\ref{tab:tap_correction} compares DC power flow accuracy on the IEEE 14-bus network with and without tap-ratio correction, measured against the MATPOWER full-AC reference.

\begin{table}[H]
\centering
\caption{DC power flow accuracy on IEEE 14-bus: tap-corrected vs.\ tap-ignored}
\label{tab:tap_correction}
\begin{tabular}{lrrr}
\toprule
\textbf{Metric} & \textbf{With tap correction} & \textbf{Without correction} & \textbf{Improvement} \\
\midrule
Max $|\Delta\theta|$ (deg) & 1.15 & 1.39 & 17.3\% \\
MAE (deg)                  & 0.57 & 0.73 & 21.2\% \\
\bottomrule
\end{tabular}
\end{table}

The tap-ratio correction reduces the mean absolute angle error by 21.2\%, with the largest improvements at buses 6, 9--14, which are electrically downstream of the corrected transformers. The residual error ($0.57^\circ$) is an inherent limitation of the lossless DC approximation, not a software defect.

% ──────────────────────────────────────────────────────────────────────────────
\section{Locational Marginal Prices}
\label{sec:lmp}
% ──────────────────────────────────────────────────────────────────────────────

Locational Marginal Prices (LMPs) represent the marginal cost of supplying one additional megawatt of load at a given bus~\cite{schweppe1988}. In the DC LOPF, they are the dual variables (shadow prices) of the nodal power balance constraints:

\begin{equation}
    \lambda_k = -\frac{\partial \text{obj}^*}{\partial P_k^\text{load}},
    \label{eq:lmp_def}
\end{equation}

where $\text{obj}^*$ is the optimal generation cost and $P_k^\text{load}$ is the load at bus $k$.

\subsection{Implementation}

LMPs are extracted directly from the JuMP model after the LOPF solve. The balance constraint at each bus $k$ is stored with a reference label, and the dual is retrieved post-solution:

\begin{lstlisting}[language=Julia, caption={LMP extraction from JuMP dual variables}, label={lst:lmp}]
# Store labelled balance constraints
balance_con = Vector{Any}(undef, n)
for k in 1:n
    balance_con[k] = @constraint(model,
        sum(B[k,m]*theta[m] for m in 1:n) == gen_inj - P_load[k] + P_shift[k])
end
optimize!(model)
# LMP = -dual(balance_con[k])
# Sign: load appears as -P_load on RHS, so d(obj)/d(P_load) = -dual
lmp = Dict(bus_names[k] => -dual(balance_con[k]) for k in 1:n)
\end{lstlisting}

\subsection{Validation}

Three scenarios verify the correctness of the LMP implementation:

\begin{enumerate}
    \item \textbf{Uncongested single-bus}: one generator at $c = 20\,\text{\euro/MWh}$ produces $\lambda = 20\,\text{\euro/MWh}$ --- the LMP equals the marginal cost.
    \item \textbf{Uncongested 3-bus}: all LMPs are equal to the marginal generator cost ($50\,\text{\euro/MWh}$ when G2 is the marginal unit), confirming no spatial price differentiation without congestion.
    \item \textbf{Congested 2-bus}: a 80\,MW line limit between a cheap area ($c_1 = 10\,\text{\euro/MWh}$) and an expensive area ($c_2 = 60\,\text{\euro/MWh}$) produces $\lambda_1 = 10$ and $\lambda_2 = 60\,\text{\euro/MWh}$, with a \emph{congestion rent} of $50\,\text{\euro/MWh}$ --- the classical result from electricity market theory.
\end{enumerate}

For the multi-period LOPF, time-varying LMPs $\lambda_k(t)$ are extracted from the per-period balance constraints, enabling analysis of how nodal prices evolve with load and renewable generation profiles.

% ──────────────────────────────────────────────────────────────────────────────
\section{Unit Commitment}
\label{sec:uc}
% ──────────────────────────────────────────────────────────────────────────────

Unit Commitment (UC) extends the LOPF by adding binary on/off decisions for generators, minimum operating constraints, and startup/shutdown costs. It is formulated as a Mixed-Integer Linear Program (MILP).

\subsection{Mathematical Formulation}

For each committable generator $g$ and time period $t$, three binary variables are introduced: $u_{g,t} \in \{0,1\}$ (commitment status), $\sigma_{g,t}^+$ (startup), and $\sigma_{g,t}^-$ (shutdown). The extended objective and key constraints are:

\begin{align}
    \min \quad & \sum_{g,t} \left[ c_g P_{g,t} + C_g^{su}\,\sigma_{g,t}^+ + C_g^{sd}\,\sigma_{g,t}^- \right]
    \label{eq:uc_obj} \\
    \text{s.t.} \quad
    & u_{g,t} - u_{g,t-1} = \sigma_{g,t}^+ - \sigma_{g,t}^-
    \label{eq:uc_transition} \\
    & p_g^\text{min} u_{g,t} \leq P_{g,t} \leq p_g^\text{max} u_{g,t}
    \label{eq:uc_dispatch} \\
    & \sum_{\tau = t - T_g^{up}+1}^{t} \sigma_{g,\tau}^+ \leq u_{g,t},
    \quad \forall t \geq T_g^{up}
    \label{eq:uc_minup} \\
    & \sum_{\tau = t - T_g^{dn}+1}^{t} \sigma_{g,\tau}^- \leq 1 - u_{g,t},
    \quad \forall t \geq T_g^{dn}
    \label{eq:uc_mindn}
\end{align}

where $C_g^{su}$, $C_g^{sd}$ are startup/shutdown costs [\euro], and $T_g^{up}$, $T_g^{dn}$ are minimum up/down times [hours].

\subsection{Implementation}

The UC solver is implemented in \texttt{julia/src/unit\_commitment.jl} and uses HiGHS as the MILP solver via JuMP. After the MILP solve, binary variables are fixed at their optimal values and a LP relaxation is re-solved to extract price-consistent LMPs --- the standard approach in electricity markets~\cite{oneill2005}. A generator is identified as committable by setting \texttt{committable=true} in the \texttt{Generator} struct, together with UC-specific fields (\texttt{min\_up\_time}, \texttt{min\_down\_time}, \texttt{startup\_cost}, \texttt{shutdown\_cost}, \texttt{initial\_status}) that map directly to the parameters of Equations~\eqref{eq:uc_logic}--\eqref{eq:uc_obj_ch2}.

\subsection{Validation}

Table~\ref{tab:uc_result} shows the 24-hour commitment schedule for a 3-bus network with a baseload unit (cheap, slow-start) and a peaker (expensive, fast-start).

\begin{table}[H]
\centering
\caption{Unit Commitment: 24-hour commitment schedule (1=ON, 0=OFF)}
\label{tab:uc_result}
\begin{tabular}{lcc}
\toprule
\textbf{Generator} & \textbf{Hours ON} & \textbf{Total starts} \\
\midrule
BaseLoad (15\,\euro/MWh, $T^{up}=4$\,h) & 9--21 (13\,h) & 1 \\
Peaker   (70\,\euro/MWh, $T^{up}=1$\,h) & 1--24 (24\,h) & 1 \\
\bottomrule
\end{tabular}
\end{table}

The MILP correctly schedules the cheap baseload unit during high-demand hours (9--21) while the more expensive peaker covers the residual. The minimum up-time constraint ($T_g^{up} = 4$\,h) is respected: once baseload starts, it operates for at least 4 consecutive hours. The average LMP across all buses is 40.2\,\euro/MWh, reflecting the weighted mix of baseload and peaker marginal costs over the 24-hour horizon.

\subsection{Unit Commitment Performance}

Table~\ref{tab:uc_benchmark_t} reports UC solve times as a function of planning horizon $T$, and Table~\ref{tab:uc_benchmark_scale} reports scaling with network size at fixed $T = 24$.

\begin{table}[H]
\centering
\caption{Unit Commitment benchmark: Julia (JuMP+HiGHS) vs PyPSA (linopy+HiGHS), 3-bus network, varying horizon, minimum time}
\label{tab:uc_benchmark_t}
\begin{tabular}{rrrr}
\toprule
\textbf{T (hours)} & \textbf{Julia (ms)} & \textbf{PyPSA (ms)} & \textbf{Speedup} \\
\midrule
  6 &  22.11 & 1{,}298 &  59$\times$ \\
 12 &  24.45 & 1{,}276 &  52$\times$ \\
 24 &  29.35 & 1{,}278 &  44$\times$ \\
\bottomrule
\end{tabular}
\end{table}

\begin{table}[H]
\centering
\caption{Unit Commitment benchmark: Julia vs PyPSA, $T=24$\,h, varying network size, minimum time}
\label{tab:uc_benchmark_scale}
\begin{tabular}{rrrr}
\toprule
\textbf{Buses} & \textbf{Julia (ms)} & \textbf{PyPSA (ms)} & \textbf{Speedup} \\
\midrule
  3 &   8.49 &   671 & 79$\times$ \\
 14 &  42.56 &   730 & 17$\times$ \\
 30 &  63.20 & 1{,}510 & 24$\times$ \\
\bottomrule
\end{tabular}
\end{table}

The UC speedup ($17$--$79\times$) is lower than for the continuous LOPF because the dominant time shifts from model construction toward MILP branch-and-bound search, which both implementations delegate to the same HiGHS solver. The PyPSA solve time is again dominated by a large fixed model-assembly overhead. The scaling with network size is non-monotonic, reflecting the combinatorial variability of the branch-and-bound search across the (synthetic) instances rather than a smooth trend.


% ──────────────────────────────────────────────────────────────────────────────
\section{AI Component: LSTM Load Forecasting and Stochastic Dispatch}
\label{sec:ai_component}
% ──────────────────────────────────────────────────────────────────────────────

The preceding sections addressed solver performance. This section introduces the AI component that justifies the "AI-assisted" framing of the thesis: an LSTM-based load forecaster with conformal prediction intervals, whose output drives a scenario-based stochastic LOPF. The implementation is contained in \texttt{julia/src/forecasting.jl} and \texttt{julia/src/stochastic\_lopf.jl}.

\subsection{Training Data}

The forecaster is trained and evaluated on \textbf{real} hourly load data for Germany, published by Open Power System Data (OPSD) from the ENTSO-E Transparency Platform \cite{opsd2020timeseries}. The cleaned series (\texttt{data/opsd\_de\_load.csv}) covers 1{,}155 complete days (2017-08-03 to 2020-09-30; peak $77{,}549$\,MW, mean $55{,}593$\,MW), peak-normalized to per-unit. A population-weighted hourly temperature series (\texttt{data/opsd\_de\_temp.csv}, six German cities, ERA5 reanalysis) is available as an optional exogenous feature. Evaluation uses a strict \emph{out-of-time} split: the LSTM is trained on the earliest 1{,}095 days and tested on the final 60 days, which it never sees during training. Because neural-network initialization is stochastic, every figure is reported as the mean $\pm$ standard deviation over three random seeds. A synthetic generator (\texttt{generate\_synthetic\_data}) is retained only as a reproducible fallback for the demonstration pipeline and unit tests.

\subsection{Forecast Quality on Real Data}
\label{sec:forecast_quality}

The honest test of a learned forecaster is whether it beats the trivial baselines on data it has never seen. Table~\ref{tab:forecast_metrics} reports the out-of-time accuracy of the LSTM against the persistence baseline (tomorrow $=$ today) and the seasonal-naive baseline (each hour $=$ the same hour one week earlier), both with and without the temperature feature.

\begin{table}[H]
\centering
\caption{Out-of-time forecast accuracy on real German load (60-day test, mean $\pm$ std over 3 seeds). Lower is better}
\label{tab:forecast_metrics}
\begin{tabular}{lccc}
\toprule
\textbf{Method} & \textbf{MAPE (\%)} & \textbf{RMSE (pu)} & \textbf{RMSE skill vs seasonal} \\
\midrule
Seasonal-naive             & \textbf{2.77} & \textbf{0.0243} & --- \\
LSTM $+$ temperature$^{*}$ & $3.67 \pm 0.68$ & $0.0300$ & $-23.4\%$ \\
LSTM (calendar only)       & $5.83 \pm 0.53$ & $0.0447$ & $-83.6\%$ \\
Persistence                & 7.73 & 0.0780 & --- \\
\bottomrule
\end{tabular}
\end{table}

\noindent$^{*}$ The temperature feature uses the \emph{actual} target-day temperature as a perfect-forecast proxy; this is an idealization, stated honestly, that upper-bounds the achievable gain from weather information.

The result is reported as it is, without spin. The LSTM \textbf{beats persistence} comfortably ($+42.7\%$ RMSE skill without temperature, $+61.5\%$ with it), but it \textbf{does not beat the seasonal-naive baseline}: with calendar features alone it is almost twice as inaccurate ($5.83\%$ versus $2.77\%$ MAPE), and even with idealized temperature it remains $23\%$ behind on average, although the best seed reaches parity ($+0.6\%$). This is a known property of aggregate national load: the weekly cycle that seasonal-naive copies for free is nearly deterministic, and a learned model adds value only through an exogenous signal --- here, weather, which is confirmed to be the decisive lever (it closes the gap from $-84\%$ to $-23\%$). The honest conclusion is that the contribution of the AI component is \emph{not} point-forecast accuracy but the calibrated uncertainty quantification developed in the next subsection, which the naive baselines cannot provide at all.

Figure~\ref{fig:ml_accuracy} visualizes the accuracy ranking; the strong seasonal-naive baseline (dashed line) sits below the LSTM in both configurations.

\begin{figure}[H]
    \centering
    \includegraphics[width=0.7\textwidth]{plots/fig_ml_accuracy.png}
    \caption{Out-of-time forecast accuracy (MAPE) on real German load. The seasonal-naive baseline is the most accurate; the LSTM beats only persistence, and temperature narrows but does not close the gap.}
    \label{fig:ml_accuracy}
\end{figure}

Although the LSTM is not the most accurate point forecaster, it provides what the naive baselines cannot: a calibrated prediction interval. Figure~\ref{fig:ml_forecast} shows one out-of-time day with the LSTM forecast, the actual load, and the 90\% conformal band. Over the full 60-day test block the empirical coverage of the band is $100\%$, comfortably above the nominal $90\%$ guarantee --- the interval is conservative because the non-conformity score is the per-day maximum absolute error, which sizes a uniform band to the worst hour.

\begin{figure}[H]
    \centering
    \includegraphics[width=0.7\textwidth]{plots/fig_ml_forecast.png}
    \caption{LSTM day-ahead forecast (weather-informed) with its 90\% conformal prediction interval, against the actual load, on an out-of-time test day. The forecast tracks the observed shape and the actual load stays within the band.}
    \label{fig:ml_forecast}
\end{figure}

\subsection{Stochastic LOPF: Propagating Forecast Uncertainty}

This subsection demonstrates the substantive contribution of the AI component: turning the calibrated forecast intervals into a risk-aware dispatch. The mechanism is illustrated on the 3-bus test network (Section~\ref{sec:test_network_definition}) with $T = 24$ hours, which keeps the optimization small enough to expose the cost distribution clearly. Three dispatch strategies are compared:

\begin{enumerate}
    \item \textbf{Static profile:} deterministic LOPF with a fixed default load profile (\texttt{DEFAULT\_LOAD\_PROFILE}) --- the simplest non-forecast baseline.
    \item \textbf{Point forecast:} deterministic LOPF with the forecaster's mean profile.
    \item \textbf{Stochastic SAA:} seven load scenarios sampled from the conformal intervals; an independent LOPF is solved per scenario and the costs are aggregated by Sample Average Approximation.
\end{enumerate}

Table~\ref{tab:stochastic_results} summarises the cost outcomes; the figures are illustrative of the mechanism on this demonstration network rather than a system-level economic claim.

\begin{table}[H]
\centering
\caption{Dispatch cost comparison: static profile vs point forecast vs stochastic LOPF on the sample out-of-time day ($T=24$\,h)}
\label{tab:stochastic_results}
\begin{tabular}{lrr}
\toprule
\textbf{Strategy} & \textbf{Total cost (\euro)} & \textbf{vs static} \\
\midrule
Static profile                    & 320{,}328 & --- \\
Point forecast (deterministic)    & 233{,}157 & $-27.2\%$ \\
Stochastic SAA --- $\mathbb{E}[\text{cost}]$ & 236{,}080 & $-26.3\%$ \\
Stochastic SAA --- CVaR$_{90}$              & 261{,}283 & $-18.4\%$ \\
Stochastic SAA --- best scenario            & 214{,}145 & $-33.1\%$ \\
Stochastic SAA --- worst scenario           & 264{,}649 & $-17.4\%$ \\
Stochastic SAA --- std dev                  & 13{,}895 & --- \\
\bottomrule
\end{tabular}
\end{table}

The point to take from Table~\ref{tab:stochastic_results} is \emph{not} the headline cost difference --- using the shaped forecast instead of a flat static profile naturally lowers cost on this day ($-27.2\%$), but that is a comparison against a deliberately weak baseline on a small demonstration network, so it is not advanced as a system-level result. The substantive output is the \textbf{risk profile} that only the stochastic formulation produces. Sampling demand scenarios across the calibrated conformal band and solving each yields an expected cost of $236{,}080$\,\euro, a worst-case (CVaR$_{90}$) of $261{,}283$\,\euro, and a spread of $\sigma = 13{,}895$\,\euro. The gap between expected and worst-case cost ($\approx 25{,}000$\,\euro, $+10.7\%$) is the reserve a risk-averse operator must budget to cover the worst $10\%$ of demand outcomes. \textbf{This risk quantification is the contribution of the AI layer, and it is inaccessible to any deterministic dispatch} --- including one driven by the more accurate seasonal-naive forecast, which returns a single number with no notion of uncertainty.

Figure~\ref{fig:scenario_costs} shows the resulting distribution of dispatch costs across the sampled scenarios, with the expected cost and the CVaR$_{90}$ marked. The spread of the distribution and the gap between the expected and worst-case cost are exactly the risk information that the stochastic formulation adds over a single deterministic solve.

\begin{figure}[H]
    \centering
    \includegraphics[width=0.7\textwidth]{plots/fig_ml_scenario_costs.png}
    \caption{Dispatch-cost distribution over demand scenarios sampled from the conformal intervals (stochastic SAA on the sample day). The solid line is the expected cost, the dashed line the CVaR$_{90}$; their gap is the tail-risk reserve.}
    \label{fig:scenario_costs}
\end{figure}

\subsection{Julia Technology: Flux.jl and Stochastic Solver}

The LSTM is implemented entirely in Julia using \textbf{Flux.jl} \cite{innes2018flux}, which provides automatic differentiation via Zygote.jl and the Adam optimizer. The full training pipeline --- data preprocessing, sequence batching, gradient computation, early stopping, and conformal calibration --- is implemented in approximately 200 lines of Julia code with no Python dependencies.

The stochastic LOPF module calls \texttt{lopf\_multiperiod} $S$ times in a loop; each call is independent, making the SAA trivially parallelisable. The current implementation is sequential (single-threaded), but parallelisation via Julia's \texttt{Threads.@threads} or distributed computing via \texttt{Distributed.jl} is a straightforward extension.

Table~\ref{tab:julia_packages_full} updates the Julia package list from Chapter~2 to include the packages added for the AI component.

\begin{table}[H]
\centering
\caption{Julia packages used in the full implementation (including AI component)}
\label{tab:julia_packages_full}
\begin{tabular}{lll}
\toprule
\textbf{Package} & \textbf{Purpose} \\
\midrule
JuMP, HiGHS   & LP/MILP modeling and solving \\
PowerModels, Ipopt & AC power flow \\
Flux.jl       & LSTM training (automatic differentiation) \\
Zygote.jl     & Reverse-mode AD backend for Flux \\
Plots.jl, StatsPlots.jl & Visualization \\
CSV, DataFrames & Benchmark result I/O \\
\bottomrule
\end{tabular}
\end{table}
